% chapters/evaluation.tex — Evaluation chapter (narrative text only)
\chapter{Evaluation Frameworks}
\label{chap:evaluation}

Rigorous evaluation of agentic \LLM{}s is challenging because agent
performance depends not only on final-answer accuracy but also on
intermediate reasoning quality, tool-use efficiency, and robustness to
distribution shift~\cite{ye2024mirai}.
\index{evaluation}
\index{benchmark}

\section{The MIRAI Benchmark}
\label{sec:eval-mirai}

MIRAI~\cite{ye2024mirai} is a benchmark for evaluating \LLM{} agents on
event forecasting: the agent must predict future geopolitical events by
reasoning over a dynamic knowledge graph constructed from news articles.
Unlike static QA benchmarks, MIRAI requires the agent to retrieve
time-sensitive information, chain multi-hop inferences, and express
calibrated uncertainty — capabilities that directly correspond to the
skills surveyed in Chapters~\ref{chap:retrieval}--\ref{chap:multimodal}.
\index{MIRAI}

\section{Dimensions of Agent Evaluation}
\label{sec:eval-dimensions}

A comprehensive agent evaluation framework should measure at least four
dimensions:

\begin{enumerate}
  \item \textbf{Task completion rate} — fraction of episodes in which the
        agent achieves the goal within the allotted steps.
  \item \textbf{Efficiency} — number of steps, tokens, and tool calls
        consumed per episode.
  \item \textbf{Robustness} — performance under observation noise,
        misleading tool outputs, or adversarial prompts.
  \item \textbf{Calibration} — correlation between the agent's expressed
        confidence and its empirical accuracy.
\end{enumerate}

\section{Connections to the Planning Objective}
\label{sec:eval-planning}

From the planning perspective of Equation~\eqref{eq:planning-objective},
evaluation reduces to measuring how closely the agent's achieved cumulative
reward approximates $\Policy^{*}$.
Current benchmarks typically proxy reward with a binary task-success
indicator, which conflates good planning with lucky execution.
Future evaluation designs should adopt denser reward signals that credit
correct intermediate reasoning steps.
